\documentclass[activite]{thedocument}
\startdatakeys
\source{originale}
\BO{2026}
\cycle{3}
\competencesDuSocle{CH2,CO2}
\motsCles{perpendicularité, distance d'un point à une droite}
\enddatakeys
\begin{document}
    \noindent{\bfseries 1) }Placer un point B qui vérifie les deux conditions~:
    \begin{itemize}
        \item \mathspoint{B} $\in$ $(d)$;
        \item la distance \mathspoint{AB} est la plus petite possible.
    \end{itemize}
    \vskip30pt{\centering
    \begin{tikzpicture}[scale=1.5]
        \coordinate[label=above left:\mathspoint{A}](A)at(0,0);
        \coordinate[label=below left:(d)](D)at(-1.5,-4);
        \draw(D)--(4,1);
        \draw node at(A){$\times$};
    \end{tikzpicture}}\par
    \noindent{\bfseries 2) } Que peut-on dire des droites $(d)$ et $(\mathspoint{AB})$~?
    \begingroup\linewidth=.95\linewidth\noindent
    \dotedline\dotedline\dotedline\endgroup
\end{document}